\scanpage{48}\mspage{35}

Prouvons $2^\circ$. On suppose donc que pour \unsure{tt} $i_0\in I$, on ait $u_{I_{<i_0},I_{/i_0}}\in C$, on veut en déduire que $u_\bullet\in\deleted{\unsure{I}}C_I$ i.e. $u_{J,I_{/i_0}}\in C$ pour \uline{tout} ouvert. \deleted{Soit $P$ l'ens \unsure{d'indices} \unsure{qui sont} \unsure{des}}

\uline{Lemme} 2.8. (Sans hypothèse artinienne sur $I$) Soit $P\subset I$ l'ens des $i_0\in I$ tels que pour \deleted{\unsure{les}} $u_\bullet\mid I_{/i_0}$ soit dans $C_{I_{/i_0}}$, i.e. que pour tout $i\le i_0$, et tout ouvert $K$ de $I_{/i}$, on ait $u_{K,I_{/i}}\in C$. \deleted{Alors} Soit $i_0\in I$ tel que $I_{<i_0}\subset P$. Alors $i_0\in P$.

\uline{Dém} Il reste à prouver que si $\deleted{\unsure{J|I}}K\subset I_{/i_0}=J$ est un ouvert, alors $u_{K,\deleted{\unsure{I_{/i_0}}}J}\in C$. OPS $K\ne I_{/i_0}$, i.e. $K\subset I_{<\unsure{i_0|i'_0}}\overset{\mathrm{def}}{=}J'$. Mais alors par hypothèse initiale sur $u$, \inserted{\unsure{on a}} $u_{J',J}\in C$, et d'autre part, \deleted{\unsure{le}} par hyp.\inserted{ sur $i_0$, \deleted{\unsure{$J$|je}}} $u_\bullet\mid J'$ est dans $C_{J'}$. Donc par le Th.\ 2.1, $1^\circ$, on a $u_{K,J'}\in C$. Donc par transitivité $u_{K,J}\in C$, qed.

\uline{Corollaire} 2.\deleted{\unsure{8|9}}\unsure{9|8}. Si $I$ est artinien, alors $P=I$, i.e. $u_\bullet\in C_I$.

Cela prouve donc 2.1. $2^\circ$. Notons que $\operatorname{Esp}(I)$ est noethérien, il en est de même de \uline{tout} ouvert (tel $J'$ ci-dessus), donc

% archive sheet number: 47

